% ---------------------------------------------------------------------
% 2.4 Potenciales relacionados con eventos y componente P300
% ---------------------------------------------------------------------
\section[POTENCIALES RELACIONADOS CON EVENTOS Y COMPONENTE P300]{Potenciales relacionados con
eventos y componente P300}
\label{sec:erp}

\subsection{Potenciales relacionados con eventos}

A diferencia de los ritmos del EEG, que se presentan de manera espontánea, los potenciales
relacionados con eventos (ERP, por sus siglas en inglés) son variaciones del voltaje que aparecen
acopladas en el tiempo a un evento específico, ya sea un estímulo sensorial, una respuesta motora o
una operación cognitiva como el reconocimiento de un estímulo o la actualización de la memoria de
trabajo \cite{kropotov2009}. La suma de los potenciales postsinápticos sincronizados con el evento
se registra en el cuero cabelludo como una serie de deflexiones positivas y negativas, conocidas
como componentes, que se nombran con la letra P o N según su polaridad y con un número que indica su
latencia aproximada en milisegundos o su posición ordinal \cite{luck2014}. La
Tabla~\ref{tab:componentes} describe las componentes que aparecen con mayor frecuencia en
paradigmas de estimulación visual.

\begin{table}[htbp]
\centering
\caption{Principales componentes de los ERP en paradigmas visuales}
\label{tab:componentes}
\small
\begin{tabularx}{\textwidth}{L{2.1cm}L{2.1cm}L{2.0cm}Y}
\toprule
\textbf{Componente} & \textbf{Latencia} & \textbf{Polaridad} & \textbf{Proceso asociado} \\
\midrule
P100 (P1) & 80--130 ms & Positiva & Procesamiento sensorial temprano en corteza visual extraestriada; se modula con la atención espacial. \\
N100 (N1) & 100--200 ms & Negativa & Procesamiento sensorial y orientación de la atención hacia el estímulo. \\
P200 (P2) & 150--250 ms & Positiva & Evaluación de las características físicas del estímulo. \\
N200 (N2) & 200--350 ms & Negativa & Detección de desajuste entre el estímulo y lo esperado; clasificación del estímulo. \\
P300 (P3) & 250--500 ms & Positiva & Actualización del contexto en la memoria de trabajo ante estímulos relevantes o poco frecuentes. \\
N400 & 300--500 ms & Negativa & Procesamiento semántico; incongruencia de significado. \\
\bottomrule
\end{tabularx}
\par\smallskip{\footnotesize Elaboración propia con información de \cite{luck2014,kropotov2009,polich2007}.}
\end{table}

\subsection{Promediado coherente}

El principal problema para observar un ERP es que su amplitud, de unos cuantos microvolts, es
considerablemente menor que la de la actividad de fondo del EEG, que para estos fines se comporta
como ruido. La forma clásica de resolverlo consiste en presentar el mismo estímulo repetidas veces,
extraer los segmentos de señal sincronizados con cada presentación y promediarlos
\cite{tong2009}. Si cada ensayo $x_i(n)$ se modela como la suma de una respuesta $s(n)$ que se
repite de forma idéntica y de un ruido $r_i(n)$ de media cero y varianza $\sigma^2$, independiente
entre ensayos, el promedio de $N$ ensayos queda como
\begin{equation}
\bar{x}(n) = \frac{1}{N}\sum_{i=1}^{N} x_i(n) = s(n) + \frac{1}{N}\sum_{i=1}^{N} r_i(n),
\label{eq:promediado}
\end{equation}
donde la desviación estándar del ruido residual se reduce a $\sigma/\sqrt{N}$, de modo que la
relación señal-ruido mejora en proporción a la raíz cuadrada del número de ensayos
\cite{luck2014}. La Figura~\ref{fig:erp} ilustra ambas ideas: del lado izquierdo, la forma de onda
típica ante un estímulo objetivo y ante uno no objetivo; del lado derecho, cómo la componente P300
va emergiendo del ruido conforme aumenta el número de ensayos promediados. Esta relación explica
por qué el P300 Speller necesita repetir varias veces la intensificación de cada fila y columna, y
por qué reducir dichas repeticiones, objetivo de este trabajo, compromete directamente la
confiabilidad de la detección.

\begin{figure}[htbp]
\centering
\includegraphics[width=\textwidth]{02-marco-referencial/erp_componentes_promediado.pdf}
\caption{(a) Forma de onda ilustrativa de un ERP ante un estímulo objetivo y uno no objetivo.
(b) Efecto del promediado de $N$ ensayos sobre la visibilidad de la componente P300. Señales
sintéticas, elaboración propia.}
\label{fig:erp}
\end{figure}

\subsection{La componente P300}

La componente P300 fue descrita por primera vez en 1965 por Sutton et al., quienes observaron que
los estímulos cuya ocurrencia era incierta para el sujeto evocaban una onda positiva de gran
amplitud alrededor de 300~ms después de su presentación \cite{sutton1965}. Años más tarde, Squires,
Squires y Hillyard distinguieron dos variantes de esta onda: la P3a, de latencia más temprana y
distribución frontocentral, provocada por estímulos novedosos que captan la atención de manera
involuntaria, y la P3b, de distribución parietal, que aparece ante estímulos relevantes para la
tarea que el sujeto está realizando \cite{squires1975}. Es esta última la que se aprovecha en el
P300 Speller, pues su aparición depende de que el usuario atienda deliberadamente al carácter que
desea seleccionar \cite{polich2007}.

La explicación más aceptada sobre el origen de la P300 es la teoría de actualización del contexto,
propuesta inicialmente por Donchin \cite{donchin1981} y posteriormente integrada por Polich
\cite{polich2007}. De acuerdo con ella, cada estímulo entrante se compara con la representación
del contexto que se mantiene en la memoria de trabajo; si no se detecta un cambio, el modelo mental
se conserva y solo se registran los potenciales sensoriales tempranos, mientras que si el estímulo
es distinto y el sujeto le asigna recursos de atención, la representación se actualiza y se genera
la P300, tal como se muestra en la Figura~\ref{fig:contexto}.

\begin{figure}[htbp]
\centering
\begin{tikzpicture}[
  node distance=0.7cm and 1.0cm,
  caja/.style={draw, rounded corners=2pt, align=center, minimum height=0.9cm, text width=2.6cm, font=\small},
  decision/.style={draw, diamond, aspect=2.1, align=center, inner sep=1pt, font=\small},
  flecha/.style={-{Stealth[length=2mm]}, thick}]
  \node[caja] (est) {Estímulo entrante};
  \node[caja, right=of est] (sens) {Procesamiento sensorial};
  \node[decision, right=of sens] (comp) {¿Mismo\\estímulo?};
  \node[caja, above right=0.3cm and 0.9cm of comp, fill=gray!12] (igual) {Se mantiene el modelo\\ \footnotesize N100, P200, N200};
  \node[caja, below right=0.3cm and 0.9cm of comp, fill=red!8] (dif) {Actualización del contexto\\ \footnotesize P300};
  \node[caja, below=0.9cm of sens, text width=3.2cm] (mem) {Memoria de trabajo\\(representación del contexto)};
  \draw[flecha] (est) -- (sens);
  \draw[flecha] (sens) -- (comp);
  \draw[thick] (comp.east) -- ++(0.35,0) coordinate (bif);
  \draw[flecha] (bif) |- node[pos=0.25, left, font=\footnotesize] {Sí} (igual);
  \draw[flecha] (bif) |- node[pos=0.25, left, font=\footnotesize] {No} (dif);
  \draw[flecha, dashed] (mem) -| (comp);
  \draw[flecha, dashed] (dif.south) |- ++(0,-0.9) -| ([xshift=0.4cm]mem.south);
\end{tikzpicture}
\caption{Modelo de actualización del contexto de la componente P300. Elaboración propia con base
en \cite{polich2007}.}
\label{fig:contexto}
\end{figure}

La P300 se caracteriza a partir de su amplitud, definida como la diferencia entre el voltaje de la
línea base previa al estímulo y el pico positivo de mayor magnitud dentro de una ventana de
búsqueda, típicamente entre 250 y 500~ms, y de su latencia, que corresponde al tiempo transcurrido
desde el inicio del estímulo hasta dicho pico \cite{polich2007}. Su distribución topográfica se
evalúa sobre los electrodos de la línea media, en los que la amplitud aumenta de las regiones
frontales hacia las parietales, alcanzando su máximo en Pz \cite{picton1992}. Tanto la amplitud
como la latencia dependen de varios factores: la amplitud es mayor mientras menos probable sea el
estímulo objetivo y más relevante resulte para la tarea, y disminuye con la fatiga, la habituación
y los intervalos cortos entre estímulos, mientras que la latencia se asocia con el tiempo que toma
clasificar el estímulo y tiende a aumentar con la edad \cite{polich2007,picton1992}.

\subsection{Paradigma del evento raro}

El procedimiento experimental más utilizado para provocar la P300 es el paradigma del evento raro,
conocido en inglés como \textit{oddball}. En él se presenta al sujeto una secuencia de estímulos de
dos categorías, una de las cuales aparece con mucha menor frecuencia que la otra, y se le pide que
atienda o cuente únicamente los estímulos de la categoría poco frecuente
\cite{farwell1988,polich2007}. Los estímulos raros, que en el contexto de una BCI corresponden al
objetivo elegido por el usuario, producen una P300 cuya amplitud es inversamente proporcional a su
probabilidad de aparición \cite{farwell1988}; la Figura~\ref{fig:oddball} representa una secuencia
típica de este paradigma.

\begin{figure}[htbp]
\centering
\includegraphics[width=0.9\textwidth]{02-marco-referencial/oddball.pdf}
\caption{Secuencia de estímulos en el paradigma del evento raro (\textit{oddball}). Elaboración
propia.}
\label{fig:oddball}
\end{figure}
